\documentclass[12pt]{article}
\usepackage[utf8]{inputenc}
\usepackage[T1]{fontenc}
\usepackage{hyperref}
\title{Mamaind\^e Language Resources for K-12 Curriculum Development}
\author{Compiled by Claude Sonnet 5.5 (claude-sonnet-5-5) for Andr\'as Kornai}
\date{}
\begin{document}
\maketitle

\section{Introduction}
Mamaind\^e is a Northern Nambikwara language of Mato Grosso, Brazil. Speaker estimates range from about 250 to 340 across four communities. The main documentation is David Eberhard's grammar. Searches turned up no school materials.

\section{Primary Educational Resources}
I found no primers, textbooks or school programmes for Mamaind\^e in search results.

\section{Linguistic and Reference Materials}
\begin{itemize}
\item David M. Eberhard, \emph{Mamaind\^e Grammar: A Northern Nambikwara Language and its Cultural Context}, LOT Dissertation Series 210, Vrije Universiteit Amsterdam, 2009 (608 pages). SIL Brazil archive record: \url{https://www.silbrazil.org/resources/archives/41276}
\item Glottolog bibliography entry for Eberhard 2009: \url{https://glottolog.org/resource/reference/id/cldf:eberhard:09}
\item Overview of the Nambikwaran languages: \url{https://en.wikipedia.org/wiki/Nambikwaran_languages}
\end{itemize}
No dictionary turned up in the searches.

\section{Religious and Community Resources}
I found no scripture translations or community organisation pages in the search results.

\section{Digital and Audio Resources}
I found no audio, video or corpus resources in the search results.

\section{Curriculum Development Context}
The language is still spoken by all adults but not by all young people, so the opportunity for intervention is now. The Eberhard grammar is a good basis for the linguistic content. A curriculum would need to be created with the community, because I saw no existing school materials.

\end{document}